\section{STaR、GRPO 与 DAPO 如何选择}

\videofigure{figures/method-selection.jpg}{不同数据规模、轨迹长度与计算预算对应不同训练方法。}{01:02:26--01:03:34}

\begin{table}[H]
\centering
\small
\begin{tabular}{>{\raggedright\arraybackslash}p{2.2cm}>{\raggedright\arraybackslash}p{3.6cm}>{\raggedright\arraybackslash}p{3.6cm}>{\raggedright\arraybackslash}p{4.0cm}}
\toprule
\textbf{方法} & \textbf{适用条件} & \textbf{主要优势} & \textbf{主要风险} \\
\midrule
STaR & 少量 rationale examples、答案可验证 & 简单高效、自举数据 & 终局过滤偏差、rationalization mismatch \\
GRPO & 可在线采样一组候选、reward 可计算 & 无 critic、系统较轻 & 全对/全错组无梯度、熵坍塌 \\
DAPO & 长 CoT、大规模在线 RL & 稳定扩展、提高有效计算率 & 实现复杂、超参耦合、监控要求高 \\
\bottomrule
\end{tabular}
\end{table}

\subsection{三个指标为何会不同}

\videofigure{figures/metrics.jpg}{RL 可能同时改变 pass@1、majority@$k$ 与 pass@$k$。}{01:03:34--01:06:20}

\begin{itemize}
    \item \textbf{pass@1}：单次采样变得更正确；
    \item \textbf{majority@$k$}：正确答案成为更常见的模态；
    \item \textbf{pass@$k$}：候选集合的多样覆盖。
\end{itemize}

RL 常显著提高 majority@$k$，却不一定提高 pass@$k$。一种解释是 policy 把概率集中到已有正确模式，提升一致性但减少多样探索；这也是“训练变强”与“覆盖变广”之间的张力。

\subsection{开放问题}

\videofigure{figures/open-problems.jpg}{Scaling RL 仍缺少对多数一致性、覆盖、多样性与自评行为的完整理论。}{01:06:20--01:10:46}

\begin{importantbox}{最值得追问的问题}
为什么 RL 让正确答案更常见，却未必让模型能解更多不同问题？回答这一问题需要同时分析 reward 形状、entropy、数据难度分布和 verifier 盲点，而不能只看平均 benchmark 分数。
\end{importantbox}

\subsection{本章小结}

STaR 适合低成本自举，GRPO 适合去 critic 的组内相对优化，DAPO 适合长 reasoning 的大规模稳定训练。方法选择应由反馈可验证性、在线采样成本和轨迹长度决定。

